\section{The Convert as Advocate}

Tertullian's datable literary career opens with a burst of apologetic. The twin
works \work{To the Nations} and the \work{Apology} are generally placed in 197
on the basis of their references to the civil wars that brought Septimius
Severus---himself an African---to sole power; the dating is a scholarly
reconstruction from internal allusions, not an ancient statement, but it is
broadly accepted. What is certain is their strategy. Instead of pleading for
pity, Tertullian indicts the legal process itself: Christians are punished for
a name, under a procedure that inverts every Roman norm, on charges---incest,
infanticide, disloyalty---that no court ever bothers to prove. Torture, used on
others to compel confession, is used on Christians to compel denial. The
argument is that of a trained forensic mind enjoying the absurdity of its
opponents' case.

Within the indictment stands a portrait of the church from inside, the first in
Latin. ``We are a body knit together as such by a common religious profession,
by unity of discipline, and by the bond of a common hope,'' he writes; the
community's fund is ``piety's deposit fund,'' spent not on banquets but ``to
support and bury poor people, to supply the wants of boys and girls destitute
of means and parents, and of old persons confined now to the house,'' and on
the shipwrecked, and on those in the mines and prisons for the Name; and the
pagans themselves testify, ``See, they say, how they love one another''
(\work{Apology} 39). Two of his sentences became permanent Christian property:
the observation, against the persecutors, that ``the oftener we are mown down
by you, the more in number we grow; the blood of Christians is seed''
(\work{Apology} 50), and the exclamation ``O noble testimony of the soul by
nature Christian!''---the soul that, startled, cries out to the one God it
pretends not to know (\work{Apology} 17).

The \work{Apology} is also the first Latin Christian book with a traceable
international readership: Eusebius cites it repeatedly through a Greek
translation (\work{Church History} 2.2.4)---the direction of translation now
reversed, Greek Christianity reading an African Latin.

Fifteen years later the advocate returned to the same brief in miniature.
\work{To Scapula}, an open letter to the proconsul Scapula who had resumed
executions, is datable with unusual precision: it mentions a portentous
darkening of the sun ``in the metropolis of Utica'' (\work{To Scapula} 3),
commonly identified with the eclipse visible in Africa on 14 August 212. The
letter contains the most remarkable sentence of early Christian political
thought: ``it is a fundamental human right, a privilege of nature, that every
man should worship according to his own convictions'' (\work{To Scapula} 2).
The claim is not modern religious liberalism---Tertullian is arguing that
coerced worship is worthless to any god---but no earlier writer is known to
have stated the principle so baldly, and it is regularly invoked in modern
discussions of religious freedom. After \work{To Scapula},
nothing Tertullian wrote can be securely dated; the advocate's last dated word
is a threat wrapped in a plea for his persecutors' own good.
